\documentclass[11pt]{article}
\usepackage[a4paper,margin=1in]{geometry}
\usepackage{amsmath,amssymb}
\usepackage[T1]{fontenc}
\usepackage{lmodern}
\usepackage{microtype}
\usepackage{xcolor}

\setlength{\parindent}{0pt}
\setlength{\parskip}{0.5em}

\newcommand{\promptbox}[1]{%
  \begingroup\setlength{\fboxsep}{10pt}\noindent
  \colorbox{blue!8}{\parbox{\dimexpr\linewidth-2\fboxsep\relax}{\textbf{Prompt. }#1}}
  \endgroup
}
\newcommand{\resultbox}[1]{%
  \begingroup\setlength{\fboxsep}{10pt}\noindent
  \colorbox{purple!8}{\parbox{\dimexpr\linewidth-2\fboxsep\relax}{\textbf{Result. }#1}}
  \endgroup
}
\newcommand{\examplebox}[1]{%
  \begingroup\setlength{\fboxsep}{10pt}\noindent
  \colorbox{green!8}{\parbox{\dimexpr\linewidth-2\fboxsep\relax}{\textbf{Example. }#1}}
  \endgroup
}

\title{\textbf{Book 1 --- Lecture 1 --- Interlude 1 --- Exercise 5}\\[2mm]
\large Orthogonal Vector Pair}
\author{}
\date{}

\begin{document}
\maketitle

\promptbox{Determine which pair of vectors are orthogonal:
\((1,1,1)\), \((2,-1,3)\), \((3,1,0)\), and \((-3,0,2)\).}

\section*{Method}
Two vectors are orthogonal exactly when their dot product is zero\footnote{For
\(\mathbf{u},\mathbf{v}\in\mathbb{R}^n\), \(\mathbf{u}\cdot\mathbf{v}=0\) means
the angle between them is \(90^\circ\).}.

\section*{Pairwise dot products}
\[
\begin{aligned}
(1,1,1)\cdot(2,-1,3) &= 4,\\
(1,1,1)\cdot(3,1,0) &= 4,\\
(1,1,1)\cdot(-3,0,2) &= -1,\\
(2,-1,3)\cdot(3,1,0) &= 5,\\
(2,-1,3)\cdot(-3,0,2) &= 0,\\
(3,1,0)\cdot(-3,0,2) &= -9.
\end{aligned}
\]

\resultbox{The only orthogonal pair is
\((2,-1,3)\) and \((-3,0,2)\).}

\examplebox{Check:
\[
(2,-1,3)\cdot(-3,0,2)=2(-3)+(-1)0+3(2)=-6+0+6=0.
\]
Since the dot product is zero, these vectors are orthogonal.}

\end{document}
